% !TeX root = ../../Book4.tex
% 纲要标题：### 8.2 Ext
% 纲要位置：计划纲要.md，第 2656 行。

\section{Ext}
\label{b4:sec:0802}

Hom 的两个变量具有相反的箭头方向，分别用投射或内射消解得到的计算必须比较。平衡性使两种模型计算同一组 Ext，并保留两变量的连接态射。


% 纲要内容（仅作写作计划，尚非教材正文）：
%
% * Hom 的右导出函子
% * 反变与协变变量
% * Ext 长正合列及具体计算
%

\subsection{Hom 的右导出函子}
\label{b4:subsec:0802:01}

\begin{definition}\label{b4:def:ext}
设 \(R\) 是含幺环，\(M,N\) 是左 \(R\) 模。对 \(N\) 取内射消解 \(N\to I^\bullet\)，定义交换群
\[
\operatorname{Ext}_R^n(M,N)\coloneqq
H^n\bigl(\operatorname{Hom}_R(M,I^\bullet)\bigr),\qquad n\ge0.
\]
这就是左正合函子 \(\operatorname{Hom}_R(M,-)\) 的右导出函子。
\end{definition}
\symidx{\(\operatorname{Ext}_R^n(M,N)\)：Ext 群}

零次是 \(\operatorname{Hom}_R(M,N)\)。在一般非交换环上，Ext 首先是交换群，不自动是左 \(R\) 模；若 \(R\) 交换，标量乘法与所有微分相容，便得到自然的 \(R\) 模结构。下标的环也不能略去，例如 \(\operatorname{Ext}_{\mathbb Z}^1(\mathbb Z/2,\mathbb Z/2)\) 与 \(\operatorname{Ext}_{\mathbb F_2}^1(\mathbb F_2,\mathbb F_2)\) 不同。

\ProseParagraphBreak
固定 \(N\) 时，\(\operatorname{Hom}_R(-,N)\) 是从左模反范畴到交换群的左正合函子。左模的投射对象正是反范畴的内射对象，所以投射消解也能给出一套右导出函子。两套计算是否一致，需要证明，不能由名字相同得出。

\subsection{反变与协变变量}
\label{b4:subsec:0802:02}

\begin{theorem}[Ext 的平衡性]\label{b4:thm:ext-balanced}
设 \(R\) 是含幺环，\(M,N\) 是幺性左 \(R\) 模，\(P_\bullet\to M\) 为投射消解，\(N\to I^\bullet\) 为内射消解。对整数 \(n\ge0\) 取
\[
\operatorname{Hom}_R(P_\bullet,N)^n=\operatorname{Hom}_R(P_n,N),
\qquad \delta^n(f)=f\,\mathrm{d}_{P,n+1}.
\]
则有自然同构
\[
H^n\bigl(\operatorname{Hom}_R(P_\bullet,N)\bigr)
\cong\operatorname{Ext}_R^n(M,N).
\]
它对 \(M\) 反变，对 \(N\) 协变，并与两变量的连接态射相容。
\end{theorem}
\begin{proof}
令 \(D^{p,q}=\operatorname{Hom}_R(P_p,I^q)\)，其中 \(p,q\ge0\)，定义
\[
\mathrm{d}_h(f)=f \mathrm{d}_P,\qquad \mathrm{d}_v(f)=(-1)^p \mathrm{d}_I f.
\]
两次水平或竖直微分为零；一水平一竖直的两种复合相差一个负号，故 \(\mathrm{d}_h \mathrm{d}_v+\mathrm{d}_v\mathrm{d}_h=0\)。总次数为 \(p+q\)，每条对角线有限，因而总复形只用有限直和。

固定 \(p\)，\(P_p\) 的投射性使 \(\operatorname{Hom}_R(P_p,-)\) 正合，所以竖列只在 \(q=0\) 有上同调 \(\operatorname{Hom}_R(P_p,N)\)。固定 \(q\)，\(I^q\) 的内射性使 \(\operatorname{Hom}_R(-,I^q)\) 正合，所以横行只在 \(p=0\) 有上同调 \(\operatorname{Hom}_R(M,I^q)\)。增广分别给出两条复形映射
\[
\operatorname{Hom}_R(P_\bullet,N)\longrightarrow\operatorname{Tot}D,
\qquad
\operatorname{Hom}_R(M,I^\bullet)\longrightarrow\operatorname{Tot}D.
\]
在第一条的 \(q=0\) 边，竖直微分为零；在第二条的 \(p=0\) 边，水平微分为零。因此它们确实与总微分相容。\cref{b4:cor:double-complex-comparison}分别沿竖列和横行说明这两条映射都是拟同构。该结果的证明只逐个消去有限对角线上的分量，不使用谱序列。

投射、内射比较映射使上述两条映射对 \(M,N\) 自然；不同选择同伦，故同调态射确定。对短正合列选择马蹄消解后，所有总复形也组成逐次短正合列；两侧增广给出它们之间的交换图。\cref{b4:thm:cohomology-long-exact}的自然性于是保证同构与连接态射相容。
\end{proof}

为了方便表示上同调类，本定理使用无符号的 \(\delta(f)=f\mathrm{d}\)。若把 \(P_n\) 放在上链次数 \(-n\)，按一般内部 Hom 规则会得到 \(\delta(f)=(-1)^{n+1}f\mathrm{d}\)；每个次数 \(n\) 乘 \((-1)^{n(n+1)/2}\) 给出这两种复形的同构。本节以后用前一种计算约定。

\ProseParagraphBreak
给定 \(u:M'\to M\)，其投射消解提升经预合成给出 \(u^*:\operatorname{Ext}_R^n(M,N)\to\operatorname{Ext}_R^n(M',N)\)；给定 \(v:N\to N'\)，后合成给出 \(v_*\)。预合成与后合成交换，所以 Ext 是两个变量的函子，而非仅对每个固定模分别定义的一组交换群。

\subsection{Ext 长正合列及具体计算}
\label{b4:subsec:0802:03}

\begin{theorem}\label{b4:thm:ext-long-exact}
设 \(R\) 是含幺环。左 \(R\) 模短正合列 \(0\to A\to B\to C\to0\) 对任意左模 \(N\) 给出自然长正合列
\[
\begin{aligned}
0&\longrightarrow\operatorname{Hom}_R(C,N)\longrightarrow\operatorname{Hom}_R(B,N)
\longrightarrow\operatorname{Hom}_R(A,N)\\
&\longrightarrow\operatorname{Ext}_R^1(C,N)\longrightarrow\operatorname{Ext}_R^1(B,N)
\longrightarrow\operatorname{Ext}_R^1(A,N)\longrightarrow\cdots.
\end{aligned}
\]
对任意左模 \(M\)，左模短正合列 \(0\to N'\to N\to N''\to0\) 给出自然长正合列
\[
\begin{aligned}
0&\longrightarrow\operatorname{Hom}_R(M,N')\longrightarrow\operatorname{Hom}_R(M,N)
\longrightarrow\operatorname{Hom}_R(M,N'')\\
&\longrightarrow\operatorname{Ext}_R^1(M,N')\longrightarrow\operatorname{Ext}_R^1(M,N)
\longrightarrow\operatorname{Ext}_R^1(M,N'')\longrightarrow\cdots.
\end{aligned}
\]
\end{theorem}
\begin{proof}
第一列取 \(N\) 的内射消解，逐次施加反变正合函子 \(\operatorname{Hom}_R(-,I^q)\)，得到 \(0\to\operatorname{Hom}(C,I)\to\operatorname{Hom}(B,I)\to\operatorname{Hom}(A,I)\to0\)，再取同调长正合列。第二列取 \(M\) 的投射消解，对给定序列逐次施加正合函子 \(\operatorname{Hom}_R(P_p,-)\)，取长正合列，再用\cref{b4:thm:ext-balanced}识别各项。增广在零次给出通常 Hom 映射，连接态射与选择独立且自然。
\end{proof}

\begin{example}\label{b4:exm:ext-cyclic}
设 \(R\) 是主理想整环，\(0\ne a\in R\)，\(N\) 是 \(R\) 模。把 \(R/(a)\) 的两项消解送入 Hom，得到 \(N\xrightarrow{a}N\)，故
\[
\operatorname{Ext}_R^0(R/(a),N)=\{\,n\in N\mid an=0\,\},
\qquad\operatorname{Ext}_R^1(R/(a),N)=N/aN,
\]
二次及以上为零。特别地，\(\operatorname{Ext}_{\mathbb Z}^1(\mathbb Z/m,\mathbb Z/n)\cong\mathbb Z/\gcd(m,n)\mathbb Z\)，其中 \(m,n>0\)。若改取 \(N=\mathbb Q\)，乘 \(m\) 为同构，零次及正次均为零；这与 \(\mathbb Q\) 内射一致。
\end{example}

在 \(R=k[\epsilon]/(\epsilon^2)\) 上，对\secref{b4:sec:0701}给出的 \(k\) 的周期消解施加 \(\operatorname{Hom}_R(-,k)\)，所有微分为零。因此 \(\operatorname{Ext}_R^n(k,k)\cong k\) 对每个 \(n\ge0\) 成立。相同的底层向量空间，在底环为 \(k\) 时正次 Ext 全为零；记录底环是计算的一部分。
